% ============================================================
% Conclusion
% ============================================================

\section{Conclusion}

Civil war does not begin in a vacuum.  Opposition groups weigh the
prospect of outside military assistance against the threat of
international resistance before they decide to fight.  I have argued
that these anticipatory calculations---the shadow of intervention---
systematically shape the decision to fight, and I have developed a theory,
a measure, and an evidentiary case for that claim.

The theoretical framework embeds the onset decision in a bargain fought in
the shadow of up to $n$ potential outside powers.  Three
assumptions---separability across interveners, non-intervention as the
reference point, and a contest--burden decomposition of the war
value---reduce the opposition's calculus to a weighted sum of
intervener-specific probabilities, moving on two axes: the net tilt of the
anticipated environment and its intensity.  Worked under the common
bargaining frictions, the model signs the axes and the government channel,
declines to sign the raw opposition channel, and fixes what the evidence can
discriminate: intensity deters in every costly-war bargain, while an
emboldening tilt requires limited commitment or divergent expectations.  The
estimates then lean against the screening account as written and toward
commitment and divergent expectations, without deciding between the two.
The framework generalizes prior work by treating both directions of
intervention, the full population of potential interveners, and their
strategic interdependence.

Operationalizing this framework required constructing a measure that did
not yet exist.  The Stage~1 classifier turns a rich directed-dyad feature
set into calibrated intervention probabilities for every potential
intervener, and multiple imputation across 25 complete datasets carries
measurement uncertainty through to the final estimates.

The onset analysis returns three findings.  First, the shadow predicts.  Added
to a structural onset model, the two aggregate shadow variables improve
out-of-sample precision under leave-one-country-out cross-validation---modestly
on their own (the area under the precision--recall curve rises from
\OosBaseAucpr\ to \OosEntAucpr) and to more than double the baseline once the
shadow is resolved by intervener type (to \OosFullAucpr)---and by out-of-sample
drop-column
importance the shadow is the model's single most valuable component, above every
structural covariate.  It also subsumes the existing proxies for anticipated
intervention: neither the Lake security hierarchy \citep{cunningham2016} nor the
\citet{gibilisco2022} structural estimates adds predictive content once the
shadow is included.  Second, the shadow enters onset along the two directions on
which the mechanisms disagree: onset falls with the overall intensity of
expectations in every measurement draw and rises with the net tilt toward the
opposition in a clear majority---the two axes the collinear channels
($r \approx \CorrEgEo$) leave identified.  The pattern leans away from the standard
common-prior screen and toward limited commitment or divergent expectations,
without deciding between them; the government-biased channel, on which contest
and burden agree, is the unambiguously deterrent one.  Third, the opposition
effect, weak as a single summed expectation, is concentrated in particular
interveners and reverses sign across them: opposition backing by a neighbor or a
rival emboldens, while backing by a major power deters---a reversal I do not
press to explain here.  The undifferentiated measure was averaging these
offsetting effects away.

The results carry implications beyond the civil war literature.  The
two-stage strategy introduced here---predict the probabilities of downstream
actions, then use those predictions as regressors in a model of the
initiating decision---applies wherever a strategic choice is made in
anticipation of others' responses.  Within the civil war literature, it
supplies what the macro-correlates framework cannot: intervention
expectations shift with the decisions of specific leaders and
governments---year to year, conflict to conflict---and so carry the temporal
variation that slow-moving structural risk factors lack.

The design also speaks to the debate over machine learning's place in
quantitative political science.  \citet{grimmerrobertsstewart2021}
distinguish measurement from causal inference and name their conflation the
cardinal sin; the two stages here keep them apart---an agnostic ensemble
evaluated on predictive criteria, then a parametric onset model carrying the
inferential weight.  \citet{knoxlucascho2022} close the loop: once a proxy
is learned, its uncertainty must be propagated and its imperfections shown
not to distort conclusions.  Under their ``average monotonicity'' condition
(\S3.1)---more of the true concept implies more of the proxy, which the
Stage~1 calibration verifies---the estimated signs of the tilt and intensity
effects are valid even if the magnitudes are attenuated.

The two stages also return opposite verdicts on flexibility: the agnostic
ensemble wins Stage~1, where tens of thousands of dyad-years and
near-deterministic exclusion zones present exactly the jagged,
interaction-heavy mapping tree-based learners are built for
\citep{montgomeryolivella2018}, while the theory-shaped logit wins Stage~2,
where onsets are rare, the index is the theory's, and flexibility buys
variance with nothing left to find.  The division is the lesson: flexibility
earns its keep where data are plentiful and the truth is jagged; structure,
where events are rare and the theory supplies the functional form.

One limitation is worth stating plainly.  The analysis rests on the
Correlates of War civil war universe, the \citet{regan2000} intervention data,
and an onset specification in the \citet{fearon2003} tradition---sources that
were standard when much of this literature was built but are now dated relative
to the UCDP/PRIO Armed Conflict Dataset \citep{gleditsch2002} that anchors
contemporary civil war research.  The ACD covers a broader conflict universe
and a more recent period, and its External Support Dataset \citep{meier2023}
codes third-party involvement at a finer, actor-level granularity than the
interstate intervention record used here.  The
measurement strategy this paper develops is not tied to these particular
sources; the estimates it produces are.

Extending the approach to that data universe is a natural next step, and the
two-stage design ports directly.  Stage~1 would train the same super-learner on
ACD onset dyads, drawing intervention labels from the External Support
Dataset---whose disaggregation by support type (troops, materiel, sanctuary,
funding) would permit a richer shadow: the anticipated \textit{forms} of
involvement, not only its bias.  The Nash fixed point and the country-year aggregation carry
over unchanged; Stage~2 would regress a contemporary onset specification on the
resulting expectations, in place of the \citet{fearon2003}-style model used
here.  Two features of the ACD would sharpen the exercise: its lower
battle-death threshold admits many smaller conflicts in which intervention is
rarer, placing a premium on rare-event metrics; and its
currency through the 2010s and 2020s would bring the Syrian, Libyan, Sahelian,
and Ukrainian intervention environments within reach.  I leave this to future
work.

Whether anticipated intervention shapes the decision to fight is a substantive
question; that it can be measured---at scale, across time, and across the full
population of potential outside powers---is the prior claim on which the
answer depends.
